% Plantilla de memoria Universidad de los Andes
% para usuarios de pdfLaTeX.
%
% Jaime Cisternas, Santiago, junio de 2004, agosto 2010
%
% Necesita todos los archivos *.tex llamados en las
% sentencias \input{} (cap1, cap2,...)
% ademas de buc.bib y las figuras.
%
% El archivo setspace.sty permite especificar el espaciado
% entre las lineas. Los archivos newapa.sty y newapa.bst
% especifican el formato de la bibliografia.
%
% Para compilar :                    pdflatex memoria
% Para generar la bibliografia :     bibtex memoria, pdflatex memoria.tex (dos veces)

\documentclass[12pt]{report}

% El archivo siguiente contiene todos los detalles
% finos del formato de la memoria.
% No es necesario modificarlo.


% Definiciones basicas de LaTeX para la memoria
%
% No modificar las lineas a continuacion a menos
% de estar muy seguro de lo que se esta haciendo!!
%
% Jaime Cisternas, julio 2004, agosto 2010

% <jcisternas@uandes.cl> 

%%%%%%%%%%%%%%%%%%%%%%%%%%%%%%%%%%%%%%%%%%%%%%%%%%%%%%%%%%%%%%%%%%%%%

% Carga librerias utiles con simbolos y estructuras
\usepackage{amsfonts}
\usepackage{amssymb}
\usepackage{amsmath}   
\usepackage{amsthm}
\usepackage{latexsym}

% Libreria para ajustar espaciado entre lineas
\usepackage{setspace}

% Libreria con estilo APA para bibliografia
% \usepackage{core/newapa}
\usepackage[backend = biber, style=apa]{biblatex}
\addbibresource{bibliografia.bib}

% Si usted está utilizando un teclado en español, puede ser útil
% usar los paquetes a continuacion. De otra manera tendrá que generar las
% tildes y la letra ñ de forma especial.
% Estos caracteres no pueden ser utilizados en etiquetas ni tampoco en modo matemático
\usepackage[utf8]{inputenc}
% \usepackage[spanish]{babel}
% Librería para obtener fecha actual
\usepackage{datetime}

% Declaraciones especificas de pdfLaTeX
\usepackage[pdftex,
        colorlinks=false,         % true or false (for final version)
        urlcolor=rltblue,         % \href{...}{...} external (URL)
	    anchorcolor=rltbrightblue,
        filecolor=weben,          % \href*{...} local file
        linkcolor=webred,         % \ref{...} and \pageref{...}
        menucolor=webdarkblue,
        citecolor=webgreen,
        pdftitle={},              % INSERT YOUR TITLE HERE
        pdfauthor={},             % INSERT YOUR NAME HERE
        pdfsubject={},
        pdfkeywords={},
        pdfpagemode=None,
        bookmarksopen=true,
	plainpages=false]{hyperref}
\usepackage[pdftex]{graphicx}
\pdfcompresslevel=9
\pdfadjustspacing=1
% Definicion de colores
\usepackage{color}
\definecolor{rltbrightred}{rgb}{1,0,0}
\definecolor{rltred}{rgb}{0.75,0,0}
\definecolor{rltdarkred}{rgb}{0.5,0,0}
\definecolor{rltbrightgreen}{rgb}{0,0.75,0}
\definecolor{rltgreen}{rgb}{0,0.5,0}
\definecolor{rltdarkgreen}{rgb}{0,0,0.25}
\definecolor{rltbrightblue}{rgb}{0,0,1}
\definecolor{rltblue}{rgb}{0,0,0.75}
\definecolor{rltdarkblue}{rgb}{0,0,0.5}
\definecolor{webred}{rgb}{0.5,.25,0}
\definecolor{webblue}{rgb}{0,0,0.75}
\definecolor{webgreen}{rgb}{0,0.5,0}
\definecolor{webdarkblue}{rgb}{0,0,0.5}
\definecolor{webbrightgreen}{rgb}{0,0.75,0}
% fin de declaraciones especificas pdfLaTeX

% tamanio de pagina
\setlength{\oddsidemargin}{.5in}
\setlength{\evensidemargin}{.0in} % en caso de usar opcion 'twoside'
\setlength{\textwidth}{6in}
\setlength{\topmargin}{-.5in}
\setlength{\textheight}{9in}

% Las siguientes definiciones pueden ser usadas en
% memorias mas matematicas
% \newtheorem{theorem}{Teorema}[section]
% \newtheorem{lemma}[theorem]{Lema}
% \newtheorem{corollary}[theorem]{Corolario}
% \newtheorem{proposition}[theorem]{Proposición}
% \newtheorem{definition}[theorem]{Definición}
% \newtheorem{claim}{Afirmación}
% \newtheorem{conjecture}[theorem]{Conjetura}
% \newtheorem{observation}[theorem]{Observación}
% \newtheorem{problem}[theorem]{Problema}

% Definicion de un ambiente para algoritmos en pseudo-codigo.
% Esta definicion puede ser mejorada.
\newtheorem{algorithm}{Algoritmo}[section]
\newcommand{\tab}{\hspace*{0.5 cm}}
% Palablas claves que deben aparecer en negrita
\newcommand{\bfWHILE}{\textbf{while~}}
\newcommand{\bfDO}{\textbf{do~}}
\newcommand{\bfEND}{\textbf{end~}}
\newcommand{\bfIF}{\textbf{if~}}
\newcommand{\bfTHEN}{\textbf{then~}}
\newcommand{\bfELSE}{\textbf{else~}}
\newcommand{\bfFOR}{\textbf{for~}}

% Nombres fijos de Latex pueden ser cambiados al Espaniol
% \renewcommand{\contentsname}{Índice General}
% \renewcommand{\chaptername}{Capítulo}
% \renewcommand{\appendixname}{Anexo}
% \renewcommand{\bibname}{Bibliografía}
% \renewcommand{\figurename}{Ilustración}
% \renewcommand{\tablename}{Tabla}
% \renewcommand{\indexname}{Índice}
% \renewcommand{\partname}{Parte}
% \renewcommand{\listfigurename}{Lista de Ilustraciones}
% \renewcommand{\listtablename}{Lista de Tablas}

% Para las enumeraciones usamos primero a,b,c,... y despues i, ii, iii,...
\renewcommand{\labelenumi}{\alph{enumi})}
\renewcommand{\labelenumii}{\roman{enumii})}
\renewcommand{\labelenumiii}{-}
\renewcommand{\labelenumiv}{-}

% Para ser consistentes en el uso de abreviaciones!
%\newcommand{\Chp}{Cap\'{\i}tulo}
%\newcommand{\Chps}{Cap\'{\i}tulos}
%\newcommand{\Sec}{Sec.} % \S
%\newcommand{\Secs}{Secs.} % \S\S
%\newcommand{\SSec}{Subsec.}
%\newcommand{\Fig}{Fig.}
%\newcommand{\Figs}{Figs.}
%\newcommand{\Eqn}{Ecn.}
%\newcommand{\Eqns}{Ecns.}

% Muestra en pantalla los nombres de los archivos usados
\listfiles

% Previene una cierta senial de ``warning: contentsline with no destination''
\newcounter{dummy}


%%%%%%%%%%%%%%%%%%%%%%%%%%%%%%%%%%%%%%%%%%%%%%%%%%%%%%%%%%%%%%%%%%%%%%%%

% LAS SIGUIENTES LINEAS DEBEN SER MODIFICADAS POR EL AUTOR DE LA MEMORIA
% SI EL TITULO ES MUY LARGO PUEDE SER CORTADO CON \linebreak

\newcommand{\nombreautor}{Luca Appiani Caro}
\newcommand{\mes}{\monthname}
\newcommand{\anio}{\the\year}
\newcommand{\titulo}{Formulation of an asset pricing model for cryptocurrencies}
\newcommand{\nombreprofuno}{Javier Mella}
\newcommand{\nombreprofdos}{Nombre Representante del Decano}
\newcommand{\nombreproftres}{Nombre Profesor Invitado}
\newcommand{\codigo}{ING-IN-001/11}

% LAS SIGUIENTES DEFINICIONES TAMBIEN SON NECESARIAS.
% ES FUNDAMENTAL QUE EL TEXTO ESTE ENCERRADO ENTRE
% PARENTESIS DE LLAVE.

\newcommand{\resumen}{
% El resumen no debe ocupar más de una página y se escribe
% al igual que el resto de la memoria en primera persona plural
% (``en este estudio hemos desarrollado...'', ``consideramos...'')
% o en voz pasiva (``el modelo ha sido analizado...'', ``será simulado...''),
% tratando de ser consistente en el uso de una de estas opciones.

% El resumen debe ser simple, fácilmente leible por cualquier persona.
% Se recomienda no incluir siglas (aún aquellas de uso habitual),
% símbolos matemáticos, ecuaciones o palabras en otro idioma.

% La estructura de la memoria ha de ser la siguiente:
% páginas de título, resumen, índice, listas de ilustraciones
% y tablas, capítulos, anexos y bibliografía.

% El formato mínimo especifica cuatro márgenes de 2.5 cm,
% tipo de letra Times de tamaño 12 puntos y un espaciado entre líneas
% parecido a 0.6 cm (1.5 veces el espaciado normal).

% Se recomienda vivamente usar esta plantilla \LaTeX{}
% ya que con ella el autor puede despreocuparse de muchos detalles
% de la edición y concentrarse en el contenido de la memoria y
% en la claridad de la exposición.
}

\newcommand{\agradecimientos}{
% Muchas gracias a todas aquellas personas que quizás
% sin saberlo y sin mala voluntad de su parte
% colaboraron a la finalización de esta investigación... 

%En esta página se puede nombrar a toda clase de personas,
%amigos y parientes, sin olvidar al profesor guía y
%a la secretaria de la Facultad.
}

\newcommand{\dedicatoria}{\it
I dedicate this work to ...
% mis papás, a mis abuelitos,
% a mi nana y a mi tía del Kinder.
}

%%%%%%%%%%%%%%%%%%%%%%%%%%%%%%%%%%%%%%%%%%%%%%%%%%%%%%%%%%%%%%%%%%%%%

% Por favor coloque todas las definiciones de simbolos a continuacion
% y no en los archivos de capitulos. Esto evitara la  existencia de
% multiples definiciones para las mismas palabras.

\newcommand{\mydef}
	{\stackrel{\mathrm{def}}{=}}
\newcommand{\e}
	{\hbox{\large{e}}}
\newcommand{\mi}
	{\hbox{\large{i}}}
\newcommand{\RR}{\mathbb{R}}
\newcommand{\dd}{\mathrm{d}}
\newcommand{\refname}{References}
%%%%%%%%%%%%%%%%%%%%%%%%%%%%%%%%%%%%%%%%%%%%%%%%%%%%%%%%%%%%%%%%%%%%%%
\begin{document}
\label{start}

% Genera las paginas de titulo, copyright, resumen,
% agradecimientos, dedicatoria, indices,...

% No es necesario modificar

% Paginas de titulo, derechos de autor y otras mas.

% Jaime Cisternas, julio 2004, agosto 2010

% declaracion especifica de pdfLaTeX
% permite incluir figuras con \includegraphics[...]{figurename}
\DeclareGraphicsExtensions{.jpg,.pdf,.mps,.png}

%%%%%%%%%%%%%%%%%%%%%%%%%%%%%%%%%%%%%%%%%%%%%%%%%%%%%%%%%%%%%%%%

  \pagenumbering{roman}
% \setcounter{page}{1}

% Escribe la pagina de titulo, incluyendo el autor y la facultad

  \thispagestyle{empty}
  \textsc{
  \vspace*{0cm}
  \begin{center}
    \Large
     Universidad de los Andes \\
     School of Engineering and Applied Sciences\\
  \end{center}
     \vspace{1cm}
  \begin{center}
     \includegraphics[angle=0,height=5cm]{logos/Logo-UANDES.png}
  \end{center}
     \vspace{1cm}
  \begin{center}
     \Large
     \titulo
  \end{center}
  \vspace{1.25cm}
  \begin{center}
    \Large
    \nombreautor
  \end{center}
  \vspace{1.25cm}
  \begin{center}
    Dissertation for the degree of \\
    Civil Industrial Engineer \\
    \vspace{1cm}
    Guiding teacher: \nombreprofuno \\
    \vspace{0.5cm}
    \codigo \\
    \vspace{0.5cm}
    Santiago, \mes\ \anio
  \end{center}
  }

%%%%%%%%%%%%%%%%%%%%%%%%%%%%%%%%%%%%%%%%%%

% Escribe la segunda pagina de titulo para ser firmada por la comision

\cleardoublepage
\thispagestyle{empty}

\begin{center}

\vspace*{2cm}
\parbox{10cm}{
\noindent
I certify that I have read this dissertation and that, in my opinion, its scope and quality are fully adequate to be considered a dissertation leading to the title of Engineer.
\vspace{1cm}

\hfill
\begin{tabular}{c}
\hspace{8cm} \\
\hline
\nombreprofuno \\
(Guiding Professor)
\end{tabular}

\vspace*{1.5cm}

\noindent
I certify that I have read this dissertation and that, in my opinion, its scope and quality are fully adequate to be considered a dissertation leading to the title of Engineer.
\vspace{0.75cm}

\hfill
\begin{tabular}{c}
\hspace{8cm} \\
\hline
\nombreprofdos
\end{tabular}

\vspace*{1.5cm}

\noindent
I certify that I have read this dissertation and that, in my opinion, its scope and quality are fully adequate to be considered a dissertation leading to the title of Engineer.
\vspace{0.75cm}

\hfill
\begin{tabular}{c}
\hspace{8cm} \\
\hline
\nombreproftres
\end{tabular}
}

\end{center}

%%%%%%%%%%%%%%%%%%%%%%%%%%%%%%%%%%%%%%%%%%

% Espaciado entre lineas
  \onehalfspacing
  %\doublespacing

% Hace pagina de derechos de autor

  \cleardoublepage
  \thispagestyle{empty}
  %\vspace*{0in}
  \begin{center}
    \copyright\ \nombreautor\ \anio \\
    All rights reserved.
  \end{center}

%%%%%%%%%%%%%%%%%%%%%%%%%%%%%%%%%%%%%%%%%%%%%%%%%%%%%%%%%%%%%%%%%%%%%%

% Genera resumen leyendo definicion de resumen

  \cleardoublepage
  \refstepcounter{dummy} \addcontentsline{toc}{section}{Summary}
  \begin{center} \Large \textbf{Summary} \end{center}
  \resumen

% Genera agradecimientos leyendo definicion de agradecimientos

  \cleardoublepage \addcontentsline{toc}{section}{Acknowledgements}
  \begin{center} \Large \textbf{Acknowledgements} \end{center}
  \agradecimientos

% Genera dedicatoria

  \cleardoublepage \vspace*{1.5in}
  \begin{flushright} \dedicatoria \end{flushright}

% Produce indice general, listas de figuras y de tablas

  \cleardoublepage
  \tableofcontents
  \cleardoublepage
  \listoffigures
  \cleardoublepage
  \listoftables
  \cleardoublepage

  \normalsize
  \pagenumbering{arabic}

%%%%%%%%%%%%%%%%%%%%%%%%%%%%%%%%%%%%%%%%%%%%%%%%%%%%%%%%%%%


%%%%%%%%%%%%%%%%%%%%%%%%%%%%%%%%%%%%%%%%%%%%%%%%%%%%%%%%%%%%%%%%%%
% cap1.tex

\chapter{Introduction}
\label{c1} % la etiqueta para referencias
\todo{Add an appropriate introduction}

\section{Literature review}
\todo{Change the references of web pages}
Due to the raising popularity in recent years of cryptocurrency, there has been a lot of research related to digital currency, from which the field of asset pricing is no exception. The latter because there is a growing interest related to the study of the factors that affect the returns of this type of assets, which certainly translates into a lot of studies whose objective is the previously mentioned. While the research topics may seem similar, it is important to note that this allows for a comprehensive categorization of the studies, despite the broadness of the related research.

The first group corresponds to the empirical studies that test for the performance of widely accepted asset pricing models such as the CAPM (\cite{sharpe1964}, \cite{Lintner1965} and \cite{mossin1966equilibrium}), FF3 \parencite{fama1993}, FF5 \parencite{fama2015}, \parencite{carhart1997}, among others. The methodology is based on the recollection of data related to returns on a specific set of cryptocurrencies in a particular time period, to then calculate the factors of the models mentioned previously. Due to the great amount of investigation that follows said framework, there are also a lot of studies that in addition to the steps mentioned previously, they complement with techniques that help understand better the underlying phenomena. For a better grasp of these groups of studies, some will be discussed that will most definitely aid the current investigation.

A first study that is included in the group of empirical studies is the one done by \parencite{gregoriou2019cryptocurrencies}. In this investigation they demonstrate that investors obtain abnormal excess returns on the London Stock Exchange from the years 2014 to 2017. The main reason behind was because of earlier studies, like \parencite{bariviera2017inefficiency}, that found evidence of inefficiency and the lack of regulation related to the cryptocurrency market. The data used corresponds to daily returns of all London Stock Exchange listed securities from the years 2014-2017, where they conclude that, applying CAPM, FF3, Carhart, and FF5, investors do indeed obtain excess returns by speculating in cryptocurrencies, suggesting that they are inefficient. While this dissertation primarily does not explore into the efficiency of cryptocurrency markets, the insights from \parencite{gregoriou2019cryptocurrencies} underline the broad applicability and versatility of such studies.

Another study is the one done by \parencite{liu2022common}, where they find that there are three factors that capture the cross-sectional expected cryptocurrency returns. Despite not forming part of the gross of the investigation, this study mentions a very interesting aspect corresponding to the different opinions people have related to cryptocurrency; they say there are two views about the related market. The first one says that all coins represent bubbles and fraud. And, the second states the technology behind said markets may become an important innovation and that at least some coins may become assets that represent stake in the future of the related technology \parencite{liu2022common}. With the current information of cryptocurrency markets it is difficult to establish right from wrong with respect to said opinions, but either way, empirical studies like \parencite{liu2022common} contribute largely to understand the factors that better explain the returns of corresponding assets. Now, regarding the research itself, the factors studied were cryptocurrency size, momentum, volume, and volatility. It is important to mention that the study focuses only on said market-factors, because financial and accounting data was not available for the cross-section of the coins that were analyzed in the data.

Regarding the conclusions drawn from said investigation, there are several to consider. Firstly, that size and momentum factors well capture the cross-section of cryptocurrnecy returns. Furthermore, a three-factor model can be constructed using market information is successful in pricing the strategies in the cryptocurrency market. A number of theoretical explanations are drawn for the factors. In relation to the cryptocurrency size premium\footnote{Smaller, less established cryptocurrencies may offer higher potential returns to compensate for their higher risk \parencite{Statista2023}.}; the cryptocurrency size factor relates to the liquidity effect\footnote{Referring to the ease with which a digital currency or token can be converted to another digital asset or cash without impacting the price and vice-versa \parencite{CFI_nodate}.}. Secondly, they find some evidence that the size premium is consistent with a mechanism proposed by cryptocurrency theories: the trade-off between capital gains and the convenience yield\footnote{According to \parencite{Investopedia2021} it is a benefit related to holding an underlying physical good, rather than the associated derivative security or contract.}.

As to momentum, the conclusions show that they are in line with the investor overreaction channel, indicating the tendency of investors to react disproportionately to new information, which in turn causes the price of cryptocurrency to swing more than it should according to its intrinsic value \parencite{10.1371/journal.pone.0264522}.

Continuing the line of empirical validation studies, \parencite{thoma2020prospect} investigated whether an investing strategy modeled by a prospect theory leads to a risk-adjusted outperformance, based on different factor models which include the \parencite{fama1993}. The utilization of prospect theory fits well in modeling the way investors inform themselves about a certain cryptocurrency, since they usually look at the price chart and then mentally represent a historical return distribution. So, according to \parencite{thoma2020prospect}, by looking at the price chart of cryptocurrency, investors evaluate the skeweness\footnote{Measure of the asymmetry of a distribution.} and evaluate the asset as a gamble, similar to lottery. The conclusions imply that cryptocurrency holders choose high prospect theory values over low values, with investors generally favoring the latter. Due to this predilection, cryptocurrencies with high prospect theory values are overbought, which reduces future gains. Cryptocurrencies with low prospect theory values on the other hand, are less likely to be overbought and might result in larger future returns. While the previous study presented did not place significant emphasis on the regression models themselves, it uses said models as a complement of the main model of the investigation, which was the prospect theory model.

Another approach to the asset pricing of cryptocurrency is the one taken by \parencite{HAYES20171308}, where a regression model is estimated using cost of production factors, rather than the usual market factors that comprise the most popular asset pricing models. Concerning the factors, \parencite{HAYES20171308} concludes that more than 84\% of value formation can be explained by three variables: computational power (as a representative for mining difficulty), rate of coin production, and the relative hardness of the mining algorithm employed.

\parencite{Shen2020} also followed a similar framework to the ones already mentioned. In this study they propose a three factor pricing model, consisting of market, size and reversal factors. Said model is compared with respect to the cryptocurrency-CAPM or C-CAPM, that uses only excess market returns to explain returns of cryptocurrency portfolios. As to the conclusions, the three-factor model based on the three factors already mentioned, has a better performance than the C-CAPM at explaining the cryptocurrency returns.

The research carried out by \parencite{Erfanian2022} provides another way of looking at the asset pricing of cryptocurrencies, while still maintaining, to some extent, the empirical studies framework mentioned in the beginning of this literature review. They apply a series of machine learning approaches to investigate whether macroeconomic, microeconomic, technical, and blockchain indicators based on economic theories can predict bitcoin price or not. Regarding the factor-based conclusions, based on a multilinear regression, the most significant long-term predictors were those of a macroeconomic nature, as well as blockchain information. 

\parencite{GROBYS20196} investigate about the popular momentum strategy implemented in the cryptocurrency market. Although there is no use of the more popular asset pricing models mentioned in the beginning, in this case a time series approach is taken, that uses the return of a security over the past months to determine the investor position on said security in the following month. In relation to the conclusions, they do not find any significant evidence as to relevant momentum payoffs in the cryptocurrency market.

Due to the extensive amount of research related to the empirical studies, a final investigation will be presented, but it is important to mention that there are much more variants of this type of studies. \parencite{Long2020} research the cross-sectional seasonality anomaly in cryptocurrency markets. Said anomaly suggests that assets with highest (lowest) average same-calendar month return tend to overperform (underperform) in the future. In simpler terms, if an investor plans to invest on a Monday, she or he should check which assets delivered the highest returns on Mondays in the past. The models used in this case include CAPM and FF3. As to the conclusions, results demonstrate that there is a strong and sizeable seasonality phenomena. However, they emphasize a limitation of their study relating the short sample period.

Now, concerning the second category of studies, they correspond to theoretic models or models that are derived from a theoretical framework. It is important to note that, unlike the empirical validation research, the quantity of theoretical models is much less. Particularly, studies focusing on cryptocurrencies are notably scarce. In despite of said shortage, one related study was found.

\parencite{koutmos2021intertemporal} developed an intertemporal regime-switching asset pricing model characterized by heterogeneous agents that have different expectations in relation to the volatility of the prices of bitcoin. The fact that the models are intertemporal, refers to the fact that the models take into account changes in market conditions and risks over time; and as to the regime-switching part, this means that said models can switch between different states or ``regimes", that could represent different market conditions. Regarding the agents, there are three: mean-variance optimizers, speculators and fundamentalists. Although the derivation of the model in this research does not come from a mathematic formulation, like the derivation of CAPM, it is interesting to review nevertheless. Through the definition of these agents, they formulate a way to represent the demand for bitcoins for each one. Then, assuming the market is only composed of said agents, they develop an asset pricing model. Finally, regarding the conclusions, one of them was that due to the special characteristics of bitcoin investors in terms of risk aversion, the fact that economic variables appear to not explain a significant part of returns is not much of a surprise. As to the models themselves, they manage to estimate the impacts of different types of investors during low and high bitcoin price volatility regimes.

Lastly, the research done by \parencite{Bennett2023}, although it does not fit into any of the two groups of studies proposed in this literature review, it provides an interesting view about different behavioral finance aspects that apply to decentralized finance\footnote{Emerging financial technology, in which cryptocurrency could be considered.}. They mention that asset pricing in rapidly evolving markets is better explained through behavioral finance, rather than through traditional finance theory. Factors like investor attention, sentiment, heuristics and biases, and network effects interact to form a highly volatile and dynamic market \parencite{Bennett2023}. A particularly compelling aspect about said research, is that presents a theoretical model of behavioral finance applications for asset pricing models related to decentralized finance, that could be taken into account when an initial proposal of factors is made. 

\subsection{Initial proposal of factors}
Having reviewed the related bibliography, determinant factors will now be proposed, which could be part of the mathematical formulation for the derivation of future models. It is important to note that the factors mentioned correspond to a preliminary proposal, and the specific way of how they could be included in the formulation of the models will not be addressed in this section.

Despite the fact that \parencite{GROBYS20196} do not find significant evidence as to relevant momentum payoffs, they evaluate only utilizing said factor as an investment strategy, but that does not mean that it does not explain the variability of cryptocurrency returns. So, in line with the conclusions outlined in the research of \parencite{liu2022common}, which does find an importance on momentum, the first factors to take into consideration are the momentum related. In order to provide greater insight, said factors are in relation to past returns (i.e. past one week returns), however the temporal aspects of said elements should be evaluated, to determine which alternative leads to better results. This type of factors could help model the behavioral finance side of the cryptocurrency market.

Furthermore, following the research of \parencite{liu2022common}, size related factors should also be studied. The inclusion of this type of variables in to a theoretic formulation could not be so straightforward, but it is important to take them into consideration because of their significance in explaining the returns of cryptocurrencies. 

Other aspects that could be accounted for correspond to representing different type of investors. In this case, the types of investors could be selected according to different characteristics, like for instance, introducing different levels of risk aversion.

Despite the fact that behavioral finance applications could be seen as endless, in terms of the different factors that could be derived from this area. Following an approach similar to the last presented research \parencite{Bennett2023} in the literature review, it would be interesting to study the viability of incorporating some factors that are of behavioral nature. Some alternatives could be: investor sentiment, investor psychological biases, or movement of other assets like commodities or stocks.

Finally, though there might be a great variety of factors that could be added to the mathematical formulation, the aggregation of them does not ensure that the model derived from said problem will explain a significant portion of the variation of the returns of cryptocurrencies. That is the reason why it is important to study whether the inclusion of a factor, significantly enhances the explanatory capacity of the model.
  
\cleardoublepage

% Bibliography
\printbibliography[heading=bibintoc, title = {References}]



\appendix % seguido de archivos de apendices

\cleardoublepage

% Escribe 'Anexos' en el Indice General
\addcontentsline{toc}{chapter}{Appendix}

% anexo_a.tex

\chapter{Technical details, tables, and others}
\label{app_a}

\section{Covariance Definition}
\label{app: cov-def}
The definition used in the derivation of the model is the one that is presented in \parencite{Intro-prob-roussas} that states the following, assuming $X$ and $Y$ are random variables that have finite expectations,
\begin{equation}
	Cov(X,Y) = E\left[(X-E(X))(Y-E(Y)\right] = E\left[XY\right] - (E\left[X\right])(E\left[Y\right])\;.
\end{equation}

\section{Stein's Lemma}
\label{app: steins lemma}
The following definition is retrieved from \parencite{balvers2001}. ``Stein's Lemma'' provides a linearization result for covariances under normality assumption when one argument is a (possible) non-linear function of a normal variable. As such it is potentially of great value in many finance applications.

It states the following:
\begin{equation}
	Cov\left[X, h\left(Y\right)\right] = E\left[h_{Y}\right]Cov\left(X,Y\right)\;,
\end{equation}
where $X$ and $Y$ are bivariate normal; $h(\cdot)$ is a differentiable function; and the subscript $Y$ indicates the partial derivative.
\section{Popular Portfolio}
\begin{table}[h!]
	\centering
	\begin{tabular}{|c|c|}
		\hline
		\textbf{Name of Cryptocurrency} & \textbf{Symbol} \\ \hline
		Bitcoin & BTC-USD \\ \hline
		Ethereum & ETH-USD \\ \hline
		Binance Coin & BNB-USD \\ \hline
		Solana & SOL-USD \\ \hline
		Ripple & XRP-USD \\ \hline
		Cardano & ADA-USD \\ \hline
		Polkadot & DOT-USD \\ \hline
		Chainlink & LINK-USD \\ \hline
		Polygon & MATIC-USD \\ \hline
	\end{tabular}
	\caption{List of Popular Cryptocurrencies portfolio and their Symbols}
	\label{tab:cryptos}
\end{table}

\section{Bull time period results}
\label{app: bull-res}

\begin{table}[h!]
	\centering
	\captionsetup{skip=0.5\baselineskip}
	\caption{Fama Mac-Beth results for the time period 14-01-2019 / 16-09-2019}
	\begin{tabular}{|c|c|}
		\hline
		\textbf{Factor} & \textbf{CAPM-model} \\ \hline
		\multicolumn{2}{|c|}{\textbf{For average returns}} \\ \hline
		Intercept & 0.0148 \\ 
		t-stat & (3.6995) \\ \hline
		Popular & -0.0001 \\ 
		t-stat & (-0.0350)\\ \hline
		Market & -0.0435 \\
		t-stat & (-1.8394)\\ \hline
		\textbf{Adjusted R-squared} & 0.4147 \\ \hline
		\multicolumn{2}{|c|}{\textbf{Cross-sec reg. for every month}} \\ \hline
		Intercept & 0.0148 \\ 
		t-stat & (1.2406) \\ \hline
		Popular & -0.0001 \\ 
		t-stat & (-0.0187)\\ \hline
		Market & -0.0435 \\
		t-stat & (-0.6275)\\ \hline
		\textbf{Average Adjusted R-squared} & 0.2861 \\ \hline
		\multicolumn{2}{|c|}{\textbf{GRS}} \\ \hline
		F-GRS (95\%) & 1.8121 \\ \hline
		F-GRS (99\%) & 2.3439 \\ \hline
		p-value & 0.1764 \\ \hline
	\end{tabular}
\end{table}

\begin{table}[h!]
	\centering
	\captionsetup{skip=0.5\baselineskip}
	\caption{Fama Mac-Beth results for the time period 07-10-2019 / 03-02-2020}
	\begin{tabular}{|c|c|}
		\hline
		\textbf{Factor} & \textbf{CAPM-model} \\ \hline
		\multicolumn{2}{|c|}{\textbf{For average returns}} \\ \hline
		Intercept & -0.0016 \\ 
		t-stat & (-0.1693) \\ \hline
		Popular & 0.0007 \\ 
		t-stat & (0.1583)\\ \hline
		Market & 0.0158 \\
		t-stat & (0.3141)\\ \hline
		\textbf{Adjusted R-squared} & -0.3109 \\ \hline
		\multicolumn{2}{|c|}{\textbf{Cross-sec reg. for every month}} \\ \hline
		Intercept & -0.0016 \\ 
		t-stat & (-0.2416) \\ \hline
		Popular & 0.0007 \\ 
		t-stat & (0.2001)\\ \hline
		Market & 0.0158 \\
		t-stat & (0.4230)\\ \hline
		\textbf{Average Adjusted R-squared} & 0.0982 \\ \hline
		\multicolumn{2}{|c|}{\textbf{GRS}} \\ \hline
		F-GRS (95\%) & 3.0122 \\ \hline
		F-GRS (99\%) & 5.0470 \\ \hline
		p-value & 0.7571 \\ \hline
	\end{tabular}
\end{table}

\begin{table}[h!]
	\centering
	\captionsetup{skip=0.5\baselineskip}
	\caption{Fama Mac-Beth results for the time period 25-02-2020 / 06-07-2020}
	\begin{tabular}{|c|c|}
		\hline
		\textbf{Factor} & \textbf{CAPM-model} \\ \hline
		\multicolumn{2}{|c|}{\textbf{For average returns}} \\ \hline
		Intercept & 0.0137 \\ 
		t-stat & (7.7126) \\ \hline
		Popular & 0.0029 \\ 
		t-stat & (2.4696)\\ \hline
		Market & -0.0207 \\
		t-stat & (-0.8020)\\ \hline
		\textbf{Adjusted R-squared} & 0.4943 \\ \hline
		\multicolumn{2}{|c|}{\textbf{Cross-sec reg. for every month}} \\ \hline
		Intercept & 0.0137 \\ 
		t-stat & (1.0953) \\ \hline
		Popular & 0.0029 \\ 
		t-stat & (0.8709)\\ \hline
		Market & -0.0207 \\
		t-stat & (-0.4435)\\ \hline
		\textbf{Average Adjusted R-squared} & 0.1186 \\ \hline
		\multicolumn{2}{|c|}{\textbf{GRS}} \\ \hline
		F-GRS (95\%) & 2.6284 \\ \hline
		F-GRS (99\%) & 4.0972 \\ \hline
		p-value & 0.0043 \\ \hline
	\end{tabular}
\end{table}

\begin{table}[h!]
	\centering
	\captionsetup{skip=0.5\baselineskip}
	\caption{Fama Mac-Beth results for the time period 31-08-2020 / 07-12-2020}
	\begin{tabular}{|c|c|}
		\hline
		\textbf{Factor} & \textbf{CAPM-model} \\ \hline
		\multicolumn{2}{|c|}{\textbf{For average returns}} \\ \hline
		Intercept & 0.0018 \\ 
		t-stat & (0.8425) \\ \hline
		Popular & 0.0002 \\ 
		t-stat & (0.1625)\\ \hline
		Market & 0.0259 \\
		t-stat & (2.4348)\\ \hline
		\textbf{Adjusted R-squared} & 0.3377 \\ \hline
		\multicolumn{2}{|c|}{\textbf{Cross-sec reg. for every month}} \\ \hline
		Intercept & 0.0018 \\ 
		t-stat & (0.2273) \\ \hline
		Popular & 0.0002 \\ 
		t-stat & (0.0532)\\ \hline
		Market & 0.0259 \\
		t-stat & (0.9777)\\ \hline
		\textbf{Average Adjusted R-squared} & 0.3462 \\ \hline
		\multicolumn{2}{|c|}{\textbf{GRS}} \\ \hline
		F-GRS (95\%) & 4.4373 \\ \hline
		F-GRS (99\%) & 9.2183 \\ \hline
		p-value & 0.9345 \\ \hline
	\end{tabular}
\end{table}

\begin{table}[h!]
	\centering
	\captionsetup{skip=0.5\baselineskip}
	\caption{Fama Mac-Beth results for the time period 24-02-2020 / 07-12-2020}
	\begin{tabular}{|c|c|}
		\hline
		\textbf{Factor} & \textbf{CAPM-model} \\ \hline
		\multicolumn{2}{|c|}{\textbf{For average returns}} \\ \hline
		Intercept & 0.0094 \\ 
		t-stat & (9.7288) \\ \hline
		Popular & 0.0030 \\ 
		t-stat & (2.4802)\\ \hline
		Market & 0.0374 \\
		t-stat & (1.7764)\\ \hline
		\textbf{Adjusted R-squared} & 0.5353 \\ \hline
		\multicolumn{2}{|c|}{\textbf{Cross-sec reg. for every month}} \\ \hline
		Intercept & 0.0094 \\ 
		t-stat & (1.3399) \\ \hline
		Popular & 0.0030 \\ 
		t-stat & (1.0963)\\ \hline
		Market & 0.0374 \\
		t-stat & (0.9218)\\ \hline
		\textbf{Average Adjusted R-squared} & 0.1895 \\ \hline
		\multicolumn{2}{|c|}{\textbf{GRS}} \\ \hline
		F-GRS (95\%) & 1.7785 \\ \hline
		F-GRS (99\%) & 2.2798 \\ \hline
		p-value & 0.3476 \\ \hline
	\end{tabular}
\end{table}

\begin{table}[h!]
	\centering
	\captionsetup{skip=0.5\baselineskip}
	\caption{Fama Mac-Beth results for the time period 31-05-2021 / 10-01-2022}
	\begin{tabular}{|c|c|}
		\hline
		\textbf{Factor} & \textbf{CAPM-model} \\ \hline
		\multicolumn{2}{|c|}{\textbf{For average returns}} \\ \hline
		Intercept & 0.0007 \\ 
		t-stat & (0.1446) \\ \hline
		Popular & 0.0042 \\ 
		t-stat & (3.7654)\\ \hline
		Market & -0.0243 \\
		t-stat & (-1.5238)\\ \hline
		\textbf{Adjusted R-squared} & 0.6039 \\ \hline
		\multicolumn{2}{|c|}{\textbf{Cross-sec reg. for every month}} \\ \hline
		Intercept & 0.0007 \\ 
		t-stat & (0.0859) \\ \hline
		Popular & 0.0042 \\ 
		t-stat & (1.0864)\\ \hline
		Market & -0.0243 \\
		t-stat & (-0.6800)\\ \hline
		\textbf{Average Adjusted R-squared} & 0.2287 \\ \hline
		\multicolumn{2}{|c|}{\textbf{GRS}} \\ \hline
		F-GRS (95\%) & 1.8740 \\ \hline
		F-GRS (99\%) & 2.4635 \\ \hline
		p-value & 0.0027 \\ \hline
	\end{tabular}
\end{table}

\begin{table}[h!]
	\centering
	\captionsetup{skip=0.5\baselineskip}
	\caption{Fama Mac-Beth results for the time period 31-01-2022 / 02-05-2022}
	\begin{tabular}{|c|c|}
		\hline
		\textbf{Factor} & \textbf{CAPM-model} \\ \hline
		\multicolumn{2}{|c|}{\textbf{For average returns}} \\ \hline
		Intercept & 0.0398 \\ 
		t-stat & (7.7552) \\ \hline
		Popular & -0.0001 \\ 
		t-stat & (-0.0919)\\ \hline
		Market & -0.0979 \\
		t-stat & (-5.7333)\\ \hline
		\textbf{Adjusted R-squared} & 0.9018 \\ \hline
		\multicolumn{2}{|c|}{\textbf{Cross-sec reg. for every month}} \\ \hline
		Intercept & 0.0398 \\ 
		t-stat & (2.1559) \\ \hline
		Popular & -0.0001 \\ 
		t-stat & (-0.0336)\\ \hline
		Market & -0.0979 \\
		t-stat & (-2.1197)\\ \hline
		\textbf{Average Adjusted R-squared} & 0.3283 \\ \hline
		\multicolumn{2}{|c|}{\textbf{GRS}} \\ \hline
		F-GRS (95\%) & 5.6931 \\ \hline
		F-GRS (99\%) & 13.6692 \\ \hline
		p-value & 0.4794 \\ \hline
	\end{tabular}
\end{table}
    % archivo de apendice

%\input{apendice_b}    % otro archivo de apendice

\label{end}
\end{document}

%%%%%%%%%%%%%%%%%%%%%%%%%%%%%%%%%%%%%%%%%%%%%%%%%%%%%%%%%%%%
